\section{Μικροϋπηρεσία Surrogate (MLP, GP, RF, k-NN)}
\label{sec:impl_surrogate}

Η υπηρεσία \texttt{surrogate\_node} υλοποιεί έναν πολυγλωσσικό πυρήνα σε Rust με εγγενείς (native) επεκτάσεις σε C++/HIP για υπολογιστική επιτάχυνση σε κάρτες γραφικών AMD (ROCm) και πολυπύρηνους επεξεργαστές (OpenMP).

\subsection{Εγγενής Πυρήνας C++/HIP και Λογική Αποστολής}
\label{subsec:impl_native_engine}

Κατά την εκκίνηση του \texttt{surrogate\_node}, εκτελείται ανίχνευση διαθέσιμου επιταχυντή:
\begin{enumerate}
    \item \textbf{ROCm / HIP (AMD GPU):} Εάν εντοπιστεί συμβατή κάρτα γραφικών AMD (π.χ. AMD Radeon RX 7600 XT, αρχιτεκτονικής RDNA3 / \texttt{gfx1102}), η εκπαίδευση του GP και ο συμπερασμός (inference) του MLP φορτώνονται στη GPU μέσω μεταγλωττισμένων πυρήνων HIP.
    \item \textbf{OpenMP (CPU):} Ως εναλλακτική λύση (fallback), ο C++ πυρήνας παραλληλοποιείται αυτόματα μέσω OpenMP σε όλους τους διαθέσιμους λογικούς πυρήνες της CPU (12 νήματα).
\end{enumerate}
Η επιλογή επιτελείται κατά την εκκίνηση μέσω δυναμικής ανίχνευσης συσκευών (\texttt{gp\_probe\_devices}) και επιλογής του βέλτιστου διαθέσιμου επιταχυντή.

\subsection{Πυρήνας Gaussian Process: Matérn 5/2}
\label{subsec:impl_gp_kernel}

Ο Gaussian Process (GP) χρησιμοποιεί πυρήνα συσχέτισης Matérn 5/2~\cite{rasmussen2006gaussian}, κατάλληλο για συναρτήσεις που είναι δύο φορές διαφορίσιμες αλλά όχι άπειρα λείες, όπως οι αεροδυναμικές επιφάνειες απόκρισης:
\begin{equation}
\label{eq:matern52}
k(r) = \left(1 + \frac{\sqrt{5}\,r}{\ell} + \frac{5r^2}{3\ell^2}\right)\exp\!\left(-\frac{\sqrt{5}\,r}{\ell}\right),
\end{equation}
όπου $r = \|\mathbf{x} - \mathbf{x}'\|$ είναι η Ευκλείδεια απόσταση μεταξύ δύο σημείων και $\ell$ το μήκος κλίμακας (length-scale) που βελτιστοποιείται με μεγιστοποίηση της οριακής πιθανοφάνειας (marginal likelihood). Το μοντέλο GP εκπαιδεύεται με βάση το σύνολο εκπαίδευσης $\mathcal{D} = \{(\mathbf{x}_i, y_i)\}_{i=1}^{N_{\text{win}}}$ και επιστρέφει τόσο την πρόβλεψη $\hat{y}$ όσο και την αβεβαιότητα $\sigma(\mathbf{x})$.

\subsection{Αρχιτεκτονική MLP}
\label{subsec:impl_mlp}

Παράλληλα με το GP, υλοποιείται ένα Multilayer Perceptron (MLP) ως μοντέλο Tier 2 χαμηλής λανθάνουσας καθυστέρησης~\cite{jin2011surrogate}. Η αρχιτεκτονική αποτελείται από:
\begin{itemize}
    \item \textbf{Είσοδος:} $d = 10$ χαρακτηριστικά (γονίδια) του διανύσματος σχεδιασμού.
    \item \textbf{Κρυφά στρώματα:} Δύο πλήρως συνδεδεμένα στρώματα διαστάσεων $64 \times 32$ νευρώνων, με ενεργοποίηση ReLU ($f(z) = \max(0, z)$).
    \item \textbf{Έξοδος:} Ένας νευρώνας γραμμικής εξόδου που προβλέπει την τιμή καταλληλότητας $\hat{F}$.
\end{itemize}
Η βελτιστοποίηση εκτελείται με τον αλγόριθμο Adam με ρυθμό μάθησης $\eta = 10^{-3}$ και ανάπτυγμα βαθμίδας (gradient clipping) σε νόρμα $1.0$, αποτρέποντας ασταθή σύγκλιση σε μικρά παράθυρα εκπαίδευσης.

\subsection{Φίλτρο Αβεβαιότητας και Ανύψωση σε Tier 3}
\label{subsec:impl_uncertainty}

Μετά από κάθε πρόβλεψη GP, ελέγχεται εάν η τυπική απόκλιση $\sigma(\mathbf{x})$ υπερβαίνει έναν εκ των προτέρων ορισμένο ουδό αξιοπιστίας $\theta_{\text{conf}}$:
\begin{equation}
\label{eq:uncertainty_filter}
\text{αν } \sigma(\mathbf{x}^*) > \theta_{\text{conf}} \implies \text{ανύψωση σε Tier 3 (πλήρης VLM).}
\end{equation}
Αυτός ο μηχανισμός εξασφαλίζει ότι γεωμετρίες που βρίσκονται σε αδύναμα καλυμμένες περιοχές του χώρου σχεδιασμού δεν αξιολογούνται μόνο από το surrogate, αλλά ανακατευθύνονται άμεσα στους εργάτες VLM για πιο αξιόπιστη εκτίμηση. Η τιμή $\theta_{\text{conf}} = 0.05$ επιλέχθηκε πειραματικά ώστε να επιτυγχάνεται ισορροπία μεταξύ ταχύτητας και ακρίβειας.

\subsection{Κυλιόμενος Ταμιευτήρας και Πολιτική Εκκαθάρισης LRU}
\label{subsec:impl_sliding_window}

Ο sliding window buffer διατηρεί τα $N_{\text{win}}$ πιο πρόσφατα έγκυρα δείγματα. Η εκκαθάριση ακολουθεί πολιτική LRU-by-generation: κάθε δείγμα σφραγίζεται με τον αριθμό γενεάς $g$ κατά την εισαγωγή του. Όταν ο ταμιευτήρας φτάσει χωρητικότητα $N_{\text{win}}$, τα δείγματα με το μικρότερο $g$ αφαιρούνται πρώτα, εξασφαλίζοντας ότι τo buffer αντανακλά πάντα τον τρέχοντα εξελικτικό χώρο και όχι απαρχαιωμένες περιοχές αναζήτησης.

\subsection{Διεπαφή gRPC: Predict και Retrain}
\label{subsec:impl_grpc_surrogate}

Η υπηρεσία εκθέτει δύο κύριες RPC μεθόδους μέσω του αρχείου πρωτοκόλλου \texttt{surrogate.proto}:
\begin{itemize}
    \item \textbf{\texttt{Predict(PredictRequest) → PredictResponse}:} Λαμβάνει το διάνυσμα γονιδίων $\mathbf{x}$ και επιστρέφει την προβλεπόμενη τιμή $\hat{F}$ μαζί με την αβεβαιότητα $\sigma$ και την ένδειξη \texttt{tier\_promoted} για ανύψωση σε Tier 3.
    \item \textbf{\texttt{Retrain(RetrainRequest) → RetrainAck}:} Ενεργοποιεί ασύγχρονη επανεκπαίδευση στο παρασκήνιο με νέα δείγματα. Η απάντηση επιστρέφεται αμέσως (\texttt{status: ACCEPTED}) χωρίς αναμονή ολοκλήρωσης, ενώ ο ενεργός δείκτης \texttt{Arc<Model>} ενημερώνεται ατομικά μόλις ολοκληρωθεί η εκπαίδευση.
\end{itemize}

